\documentclass{beamer}

% Theme selection
\usetheme{Madrid}
\usecolortheme{default}

% Packages
\usepackage[utf8]{inputenc}
\usepackage{graphicx}
\usepackage{amsmath}
\usepackage{amssymb}
\usepackage{hyperref}

% Title page information
\title[S2S Translation]{Speech to Speech Translation Systems}
\subtitle{Benchmarking and Implementation}
\author[Tanay, Jay, Kshitij, Vatsal]{Tanay Bhat  (22B3303) \\ Jay Mehta (22B1281) \\ Kshitij Vaidya (22B1829) \\ Vatsal Melwani (22B0396)}
\institute[IITB]{IIT Bombay}
\date{\today}

\begin{document}

% Title slide
\begin{frame}
\titlepage
\end{frame}

% Table of contents
\begin{frame}{Outline}
\tableofcontents
\end{frame}

% Introduction Section
\section{Introduction}

\begin{frame}{Introduction}

\end{frame}

% ASR Section
\section{Automatic Speech Recognition (ASR)}

\begin{frame}{ASR: Overview}
The ASR module of a S2S pipeline first converts the input speech into text. It involves several key components:
\begin{itemize}
    \item \textbf{Audio Preprocessing:} Noise reduction and feature extraction using MFCCs, Mel spectrograms etc.
    \item \textbf{Encoder:} Responsible for capturing temporal dependencies and contextual information of the input audio using the extracted features. Typically implemented using RNNs, CNNs, or Transformers.
    \item \textbf{Decoder:} Generates the output text sequence based on the encoder reprensentation. Current state-of-the-art models often use attention mechanisms to focus on relevant parts of the input during decoding.
\end{itemize}
\end{frame}

\begin{frame}{ASR Benchmarking: Setup}
\begin{itemize}
\item Accent and dialect variation
\item Background noise handling
\item Real-time processing requirements
\item Domain-specific vocabulary
\end{itemize}
\end{frame}

\begin{frame}{ASR Benchmarking: Results}
\begin{itemize}
\item Accent and dialect variation
\item Background noise handling
\item Real-time processing requirements
\item Domain-specific vocabulary
\end{itemize}
\end{frame}

\begin{frame}{ASR Benchmarking: Inferences}
\begin{itemize}
\item Accent and dialect variation
\item Background noise handling
\item Real-time processing requirements
\item Domain-specific vocabulary
\end{itemize}
\end{frame}

% MT Section
\section{Machine Translation (MT)}

\begin{frame}{MT: Overview}
\begin{itemize}
\item Translates text from source to target language
\item Evolution: Rule-based → Statistical → Neural
\item Transformer architecture revolution
\item Quality assessment and evaluation
\end{itemize}
\end{frame}

\begin{frame}{MT: Key Approaches}
\begin{itemize}
\item Neural Machine Translation (NMT)
\item Transfer learning and multilingual models
\item Low-resource language handling
\item Context-aware translation
\end{itemize}
\end{frame}

% TTS Section
\section{Text-to-Speech (TTS)}

\begin{frame}{TTS: Overview}
\begin{itemize}
\item Converts text into natural-sounding speech
\item Components: Text analysis, prosody generation, synthesis
\item Neural TTS models (WaveNet, Tacotron)
\item Voice cloning and personalization
\end{itemize}
\end{frame}

\begin{frame}{TTS: Quality Factors}
\begin{itemize}
\item Naturalness and intelligibility
\item Prosody and emotion
\item Multi-speaker capabilities
\item Real-time generation
\end{itemize}
\end{frame}

% Full Pipeline Section
\section{Full Pipeline Integration}

\begin{frame}{End-to-End Pipeline}
\begin{itemize}
\item ASR → MT → TTS workflow
\item Use case: Real-time speech translation
\item Latency and quality trade-offs
\item Error propagation considerations
\end{itemize}
\end{frame}

\begin{frame}{Pipeline Architecture}
\begin{itemize}
\item Modular vs. end-to-end approaches
\item API integration strategies
\item Scalability and deployment
\item Performance optimization
\end{itemize}
\end{frame}

\begin{frame}{Applications}
\begin{itemize}
\item International conference interpretation
\item Language learning platforms
\item Accessibility tools
\item Customer service automation
\end{itemize}
\end{frame}

% Conclusion Section
\section{Conclusion}

\begin{frame}{Conclusion}
\begin{itemize}
\item Integrated ASR-MT-TTS systems enable powerful applications
\item Continued advances in neural architectures
\item Challenges remain in low-resource scenarios
\item Future directions: multimodal integration, personalization
\end{itemize}
\end{frame}

\begin{frame}{Thank You}
\begin{center}
{\Large Thank you for your attention!}

\vspace{1cm}

Questions?

\vspace{1cm}

\small
Contact: your.email@example.com
\end{center}
\end{frame}

\end{document}